\lecture{8}
\title{Undecidability}

\begin{document}

During the first half hour of class,
we go over the proofs of $|\mathbb{N}| = |\mathbb{Q}|$
and $|\mathbb{N}| \neq |\mathbb{R}|$ as a warm-up.

\subsection{$A_{\mathrm{TM}}$ Undecidability}

\begin{theorem}{$A_{\mathrm{TM}}$}
    Consider the language
    $A_{\mathrm{TM}} := \{\langle M, w \rangle \mid
    M \text{ is a TM that accepts } w\}$.
    Then $A_{\mathrm{TM}}$ is not decidable.
\end{theorem}

\begin{proof}
    Suppose FTSOC that some TM $H$ decides $A_{\mathrm{TM}}$.
    Recall that this means that:
    \[
    H \text{ on } \langle M, w \rangle = \begin{cases}
    \textsc{Accept} & \text{ if } M \text{ accepts } w. \\
    \textsc{Reject} ~ \textit{ (via halting)} & \text{ if } M \text{ rejects } w.
    ~ \textit{ (via halting or looping)}
    \end{cases}
    \]
    Now consider the TM $D$ that takes
    an encoding of a Turing Machine $\langle M \rangle$ as input
    and does the following:
    \[
    D \text{ on } \langle M \rangle
    = \begin{cases}
    \textsc{Accept} & \text{ if } H \text{ rejects on }
    \langle M, \langle M \rangle \rangle. \\
    \textsc{Reject} & \text{ if } H \text{ accepts on }
    \langle M, \langle M \rangle \rangle.
    \end{cases}
    ~ \iff ~
    \text{``}D \text{ accepts } \langle M \rangle
    \text{ iff } M \text{ does not accept } \langle M \rangle\text{.''}
    \]
    The contradiction, now, is that the definition above implies
    that $D$ accepts $\langle D \rangle$ if and only if
    $D$ rejects $\langle D \rangle$. ~ $\blacksquare$
\end{proof}

To see how the argument above relates to diagonalization,
consider the diagram below.
\begin{center}
    \frame{\includegraphics[width=0.35\textwidth]{lecture-08/ATM-diagonalization}}
\end{center}
The entry in row $M_i$ and column $\langle M_j \rangle$
indicates whether $M_i$ accepts or rejects on input $\langle M_j \rangle$.

\subsection{Turing-Unrecognizable}

We first require the following lemma.

\begin{theorem}{Decidable vs. Turing-Recognizable}
    A language $A$ is decidable
    if and only if $A$ and $\overline{A}$
    are both Turing-recognizable.
\end{theorem}

\begin{proof}
    The $\implies$ direction is easy;
    combine the following two observations to win.
    \begin{itemize}
        \item In general, if $A$ is decidable,
        then its complement $\overline{A}$ is also decidable---just
        flip the answer.
        \item In general, if $A$ is decidable,
        then $A$ is also Turing-recognizable.
    \end{itemize}
    Now for the $\impliedby$ direction.
    Suppose $R$ recognizes $A$ and $S$ recognizes $\overline{A}$.
    Note for any input $w$ the following:
    \[
    \text{``If $w \in A$, then $R$ must eventually accept.''}
    ~~~ \parallel ~~~
    \text{``If $w \in \overline{A}$, then $S$ must eventually accept.''}
    \]
    So we run $R$ and $S$ on $w$ in parallel
    until one of them accepts;
    return $\textsc{Accept}$ if $R$ accepts,
    and return $\textsc{Reject}$ if $S$ accepts. ~ $\blacksquare$
\end{proof}

We can now use this lemma to prove that languages are not Turing-recognizable.

\begin{theorem}{$\overline{A_{\mathrm{TM}}}$}
    The language
    $\overline{A_{\mathrm{TM}}} := \{\langle M, w \rangle \mid
    M \text{ is a TM that does not accept } w\}$
    is not Turing-recognizable.
\end{theorem}

\begin{proof}
    Recall the following:
    \begin{itemize}
        \item From the end of
        \href{/18.404/lectures/7/#algorithms-on-grammars-and-turing-machines}{Lecture 7},
        we know that $A_{\mathrm{TM}}$ is Turing-recognizable.
        \item From earlier in this lecture,
        we know that $A_{\mathrm{TM}}$ is \emph{not} decidable.
    \end{itemize}
    The previous theorem now implies
    that $\overline{A_{\mathrm{TM}}}$ cannot be Turing-recognizable.
    ~ $\blacksquare$
\end{proof}

\end{document}